% !TEX program = lualatex
% Coverage example: \mfr called explicitly by the author. See #90.
%
% \mfr is public API -- the README's Memorandum lines table lists it first -- and until
% this file nothing in the tree called it. Its only occurrence was inside a % comment at
% examples/example-grid.tex:34.
%
% This is NOT redundant with example-nomemoline.tex. That file reaches \mfr through
% \am@memoline's default branch, which runs at \AtEndDocument. This one calls \mfr from
% the preamble, so the \listadd happens before \begin{document} rather than during the
% render. Both routes must work, and only the timing distinguishes them; a change that
% made \mfr safe in one context and not the other would pass with a single example.
%
% Rendered output is expected to be identical to example-nomemoline's, which is the point:
% declaring nothing and declaring \mfr mean the same thing.
\documentclass{../armymemo}

\address{Organizational Name/Title}
\address{Standardized Street Address}
\address{CITY, STATE~~12345-1234}

\author{John W. Smith}\rank{CPT}\branch{CY}
\officesymbol{ABC-DEF-GH}
\signaturedate{10 April 2019}
\mfr
\documentmark{UNCLASSIFIED//FOR OFFICIAL USE ONLY (EXAMPLE ONLY)}
\subject{A memorandum declaring its line with the record shorthand}
\title{COMMANDING}

\begin{document}

\begin{enumerate}
\item This memorandum calls \verb!\mfr! in its preamble. The class renders MEMORANDUM FOR
  RECORD from the list, not from the default branch.
\item Compare example-nomemoline.tex, which declares nothing and reaches the same text
  through \verb!\am@memoline!'s default. The two bodies are worded identically so that any
  divergence between the routes shows up as a difference between the goldens.
\end{enumerate}

\end{document}
